\documentclass[a4paper,11pt]{article}
\usepackage[T2A]{fontenc}
\usepackage[utf8]{inputenc}
\usepackage[russian]{babel}
\usepackage{amsmath,amssymb,mathtools}
\renewcommand{\familydefault}{\sfdefault}
\usepackage{icomma}
\usepackage[a4paper,margin=20mm,headheight=26pt,headsep=8mm,footskip=12mm]{geometry}
\usepackage{microtype}
\usepackage{needspace}
\usepackage{tikz}
\usetikzlibrary{arrows.meta,positioning,shapes.geometric,calc}
\usepackage{tabularx,booktabs,array}
\usepackage{enumitem}
\usepackage{hyperref}
\usepackage{parskip}
\usepackage{fancyhdr}
\usepackage{setspace}
\usepackage{xcolor}

\definecolor{accent}{HTML}{2C6E6F}
\definecolor{darkaccent}{HTML}{214F50}
\definecolor{textgray}{HTML}{2F3333}
\definecolor{softgray}{HTML}{687475}
\definecolor{linegray}{HTML}{D5DDDD}
\hypersetup{colorlinks=true,linkcolor=darkaccent,urlcolor=darkaccent}
\setlength{\parindent}{0pt}
\setstretch{1.16}
\setlist[itemize]{leftmargin=1.5em,topsep=0.2em,itemsep=0.22em}
\setlist[enumerate]{leftmargin=1.6em,topsep=0.2em,itemsep=0.65em}
\pagestyle{fancy}
\fancyhf{}
\fancyhead[L]{\small\color{softgray} Действия со смешанными числами}
\fancyhead[R]{\small\color{softgray} 10.10.26}
\fancyfoot[C]{\small\color{softgray}\thepage}
\renewcommand{\headrulewidth}{0.3pt}
\renewcommand{\headrule}{\hbox to\headwidth{\color{linegray}\leaders\hrule height \headrulewidth\hfill}}

\newcommand{\sectiontitle}[1]{%
  \Needspace{6\baselineskip}%
  \vspace{0.34em}{\large\bfseries\color{darkaccent}#1}\par
  \vspace{0.16em}{\color{linegray}\hrule}\vspace{0.42em}}
\newcommand{\idea}[1]{\textbf{Ключевая идея.} #1}
\newcommand{\note}[1]{\textbf{Обратите внимание.} #1}
\newcommand{\example}[1]{\textbf{Пример.} #1}
\newcommand{\mixed}[3]{#1\,\frac{#2}{#3}}

\begin{document}
\setlength{\abovedisplayskip}{8pt}
\setlength{\belowdisplayskip}{8pt}
\setlength{\abovedisplayshortskip}{6pt}
\setlength{\belowdisplayshortskip}{6pt}
\color{textgray}
\begin{titlepage}
\thispagestyle{empty}
\centering
\vspace*{27mm}
{\small\color{softgray} МАТЕМАТИКА}\par
\vspace{9mm}
{\Huge\bfseries\color{darkaccent} Чек-лист}\par
\vspace{8mm}
{\LARGE\bfseries Действия со смешанными числами}\par
\vspace{3mm}
{\large сложение, вычитание, скобки, уравнения и задачи}\par
\vspace{17mm}
{\large 10.10.26}\par
\vfill
{\large\bfseries Лёвин Артём Александрович}\par
\vspace{2mm}
{\normalsize эксклюзивно для Нади Клименко}\par
\vspace{19mm}
\begin{tikzpicture}[remember picture,overlay]
  \draw[accent,line width=1.05pt]
  ($(current page.south west)+(18mm,20mm)$)
  .. controls ($(current page.south west)+(65mm,30mm)$)
  and ($(current page.south east)+(-66mm,10mm)$)
  .. ($(current page.south east)+(-18mm,20mm)$);
\end{tikzpicture}
\end{titlepage}

\sectiontitle{1. Как устроено смешанное число}
\idea{Смешанное число состоит из целой и дробной частей. С целыми частями и дробными частями удобно работать отдельно.}
\[
\mixed{a}{b}{c}=a+\frac bc,\qquad 0\le b<c,\quad c>0.
\]
Например, \(\mixed{4}{5}{12}=4+\frac5{12}\). При сложении и вычитании дробные части сначала приводят к общему знаменателю.

\sectiontitle{2. Сложение смешанных чисел}
\textbf{Алгоритм:} приведите дробные части к общему знаменателю; сложите целые части и числители; если дробная часть неправильная, выделите целую; сократите дробь.

\example{Вычислим \(\mixed{14}{5}{12}+\mixed{17}{1}{6}\).}
\[
\begin{aligned}
\mixed{14}{5}{12}+\mixed{17}{1}{6}
&=\mixed{14}{5}{12}+\mixed{17}{2}{12}\\
&=(14+17)+\frac{5+2}{12}
=\boxed{\mixed{31}{7}{12}}.
\end{aligned}
\]

Если сумма дробных частей больше единицы, её нужно преобразовать:
\[
\mixed{3}{7}{10}+\mixed{2}{9}{10}
=5+\frac{16}{10}=6+\frac6{10}
=\boxed{\mixed{6}{3}{5}}.
\]

\sectiontitle{3. Вычитание: когда занимать единицу не требуется}
Сначала приведите дробные части к общему знаменателю, затем вычтите целые части и дробные части.

\example{Вычислим \(\mixed{4}{5}{6}-\mixed{2}{1}{4}\).}
\[
\begin{aligned}
\mixed{4}{5}{6}-\mixed{2}{1}{4}
&=\mixed{4}{10}{12}-\mixed{2}{3}{12}\\
&=(4-2)+\frac{10-3}{12}
=\boxed{\mixed{2}{7}{12}}.
\end{aligned}
\]

\sectiontitle{4. Вычитание с «заниманием» единицы}
\idea{Если после приведения к общему знаменателю дробная часть уменьшаемого меньше дробной части вычитаемого, возьмите одну единицу из целой части уменьшаемого.}
\[
\mixed{a}{b}{c}=(a-1)+\frac{b+c}{c}\qquad (a\ge1).
\]
\example{Вычислим \(\mixed{5}{7}{12}-\mixed{1}{5}{8}\).}
\[
\begin{aligned}
\mixed{5}{7}{12}-\mixed{1}{5}{8}
&=\mixed{5}{14}{24}-\mixed{1}{15}{24}\\
&=\mixed{4}{38}{24}-\mixed{1}{15}{24}\\
&=(4-1)+\frac{38-15}{24}
=\boxed{\mixed{3}{23}{24}}.
\end{aligned}
\]
\note{При «занимании» целая часть уменьшается на единицу, а числитель увеличивается ровно на знаменатель: \(14+24=38\).}

\sectiontitle{5. Если из целого числа вычитают дробь}
Представьте одну единицу в виде дроби с нужным знаменателем:
\[
1-\frac45=\frac55-\frac45=\boxed{\frac15},
\qquad
2-\frac37=1+\frac77-\frac37=\boxed{\mixed{1}{4}{7}}.
\]
Если вычитается смешанное число, сначала учитывайте обе его части:
\[
6-\mixed{1}{5}{9}=5+\frac99-\mixed{1}{5}{9}
=\boxed{\mixed{4}{4}{9}}.
\]

\sectiontitle{6. Скобки и несколько действий}
\textbf{Удобный порядок:} вычислите доступные выражения внутри скобок; затем приведите оставшиеся дроби к общему знаменателю; при необходимости займите единицу.

\example{Вычислим.}
\[
\begin{aligned}
\frac16-\left(1-\frac{17}{18}\right)
&=\frac16-\frac1{18}\\
&=\frac3{18}-\frac1{18}
=\frac2{18}=\boxed{\frac19}.
\end{aligned}
\]

\example{Вычтем последовательно два смешанных числа.}
\[
\begin{aligned}
\mixed{6}{1}{6}-\mixed{2}{5}{24}-\mixed{3}{11}{12}
&=\mixed{6}{8}{48}-\mixed{2}{10}{48}-\mixed{3}{44}{48}\\
&=\mixed{5}{56}{48}-\mixed{2}{10}{48}-\mixed{3}{44}{48}\\
&=(5-2-3)+\frac{56-10-44}{48}
=\boxed{\frac1{24}}.
\end{aligned}
\]
\note{Если в выражении несколько вычитаемых, сравните дробную часть уменьшаемого с \emph{суммой} вычитаемых дробных частей.}

\sectiontitle{7. Смешанные и десятичные дроби в одном выражении}
Переведите десятичную дробь в обыкновенную или смешанную, затем примените знакомые действия:
\[
3{,}4=\frac{34}{10}=\frac{17}{5}=\mixed{3}{2}{5}.
\]
\[
\begin{aligned}
\mixed{2}{3}{4}+3{,}4
&=\mixed{2}{3}{4}+\mixed{3}{2}{5}\\
&=5+\frac{15+8}{20}
=5+\frac{23}{20}
=\boxed{\mixed{6}{3}{20}}.
\end{aligned}
\]
\note{В конце выразите ответ в требуемом виде и сократите дробную часть.}

\sectiontitle{8. Простейшие уравнения со смешанными числами}
\idea{Для уравнения \(x+a=b\) найдите неизвестное слагаемое: \(x=b-a\). При вычислении разности используйте «занимание» единицы.}

\example{Решим \(x+\mixed{3}{8}{13}=6\).}
\[
\begin{aligned}
x&=6-\mixed{3}{8}{13}\\
 &=\mixed{5}{13}{13}-\mixed{3}{8}{13}
 =\boxed{\mixed{2}{5}{13}}.
\end{aligned}
\]
\textbf{Проверка:} \(\mixed{2}{5}{13}+\mixed{3}{8}{13}=6\).

\sectiontitle{9. Текстовая задача на совместную работу}
\example{Первая труба заполняет резервуар за 6 часов, вторая --- за 8 часов. Какую часть резервуара останется заполнить после одного часа их совместной работы?}

Весь объём резервуара принимаем за \(1\). Производительность --- часть резервуара за час.
\begin{center}
\renewcommand{\arraystretch}{1.35}
\begin{tabularx}{\linewidth}{>{\raggedright\arraybackslash}Xcc}
\toprule
Работа & Время на весь объём & За один час\\
\midrule
Первая труба & \(6\) ч & \(\frac16\)\\
Вторая труба & \(8\) ч & \(\frac18\)\\
Обе вместе & --- & \(\frac16+\frac18=\frac7{24}\)\\
\bottomrule
\end{tabularx}
\end{center}
За час обе трубы заполняют \(\frac7{24}\) резервуара. Следовательно, останется:
\[
1-\left(\frac16+\frac18\right)
=\frac{24}{24}-\frac7{24}
=\boxed{\frac{17}{24}}.
\]
\note{В вопросах «сколько останется» из целого \(1\) вычитайте уже выполненную часть работы.}

\sectiontitle{10. Самопроверка перед тренировкой}
\begin{itemize}[label={$\square$}]
  \item Привожу дробные части к одному знаменателю до сложения или вычитания.
  \item При вычитании проверяю, достаточно ли дробной части уменьшаемого.
  \item Если занимаю единицу, уменьшаю целую часть ровно на \(1\) и добавляю \(\frac cc\).
  \item Учитываю скобки и порядок действий.
  \item Десятичные дроби перевожу в обыкновенные без потери целой части.
  \item В текстовой задаче отличаю выполненную часть работы от оставшейся.
  \item Выделяю целую часть из неправильной дроби и сокращаю результат.
\end{itemize}

\clearpage
\begin{center}
{\Large\bfseries\color{darkaccent} Алгоритм вычитания смешанных чисел}
\end{center}
\vspace{7mm}
\begin{center}
\begin{tikzpicture}[
node distance=12mm,
start/.style={ellipse,draw=accent,line width=.9pt,text width=58mm,minimum height=12mm,align=center,font=\small},
process/.style={rounded corners=2mm,draw=accent,line width=.9pt,text width=105mm,minimum height=15mm,align=center,font=\small,inner sep=3mm},
decision/.style={diamond,aspect=2.7,draw=accent,line width=.9pt,text width=48mm,align=center,font=\small,inner sep=1.4mm},
arrow/.style={-{Stealth[length=2.5mm]},line width=.85pt,color=darkaccent}
]
\node[start] (begin) {Даны два смешанных числа};
\node[process,below=of begin] (common) {Приведите дробные части к общему знаменателю};
\node[decision,below=15mm of common] (compare) {Дробная часть уменьшаемого меньше?};
\node[process,below left=19mm and 12mm of compare,text width=44mm] (borrow) {Займите единицу: целую часть уменьшите на 1, к числителю прибавьте знаменатель};
\node[process,below right=19mm and 12mm of compare,text width=44mm] (no) {Оставьте целую и дробную части без изменений};
\node[process,below=37mm of compare] (subtract) {Вычтите целые части и дробные части по отдельности};
\node[process,below=of subtract] (finish) {Если нужно, выделите целую часть и сократите дробь};
\node[start,below=of finish] (end) {Ответ};
\draw[arrow] (begin) -- (common);
\draw[arrow] (common) -- (compare);
\draw[arrow] (compare.west) -| node[pos=.2,above,font=\small] {да} (borrow.north);
\draw[arrow] (compare.east) -| node[pos=.2,above,font=\small] {нет} (no.north);
\draw[arrow] (borrow.south) |- (subtract.west);
\draw[arrow] (no.south) |- (subtract.east);
\draw[arrow] (subtract) -- (finish);
\draw[arrow] (finish) -- (end);
\end{tikzpicture}
\end{center}

\clearpage
\sectiontitle{Тренировочный блок --- Путь А}
\textbf{10 заданий.} Решайте по порядку; записывайте основные действия и сокращайте ответы. В заданиях с десятичной дробью сначала переведите её в обыкновенную.
\begin{enumerate}
 \item Вычислите: \(\mixed{2}{1}{5}+\mixed{3}{2}{5}\).
 \item Вычислите: \(\mixed{6}{5}{12}+\mixed{1}{3}{8}\).
 \item Вычислите: \(\mixed{7}{5}{6}-\mixed{4}{1}{3}\).
 \item Вычислите: \(\mixed{5}{1}{6}-\mixed{2}{3}{4}\).
 \item Вычислите: \(6-\mixed{1}{5}{9}\).
 \item Вычислите: \(\frac16-\left(1-\frac{17}{18}\right)\).
 \item Вычислите: \(\mixed{2}{3}{4}+3{,}4\).
 \item Вычислите:
 \[
 \mixed{6}{1}{6}-\mixed{2}{5}{24}-\mixed{3}{11}{12}.
 \]
 \item Решите уравнение: \(x+\mixed{3}{7}{12}=\mixed{7}{1}{8}\).
 \item Первая труба заполняет резервуар за \(5\) часов, вторая --- за \(12\) часов. Какую часть резервуара останется заполнить после \(2\) часов их совместной работы?
\end{enumerate}

\clearpage
\sectiontitle{Ответы}
\renewcommand{\arraystretch}{1.48}
\begin{tabularx}{\linewidth}{>{\bfseries}c X}
\toprule
№ & Ответ\\
\midrule
1 & \(\mixed{5}{3}{5}\)\\
2 & \(\mixed{7}{19}{24}\)\\
3 & \(\mixed{3}{1}{2}\)\\
4 & \(\mixed{2}{5}{12}\)\\
5 & \(\mixed{4}{4}{9}\)\\
6 & \(\frac19\)\\
7 & \(\mixed{6}{3}{20}\)\\
8 & \(\frac1{24}\)\\
9 & \(\mixed{3}{13}{24}\)\\
10 & \(\frac{13}{30}\) резервуара\\
\bottomrule
\end{tabularx}

\vspace{8mm}
\textbf{Проверка трудных мест.}
\begin{itemize}
 \item В № 4, 5 и 8 нужно занимать единицу.
 \item В № 7 десятичная дробь \(3{,}4=\mixed{3}{2}{5}\).
 \item В № 9 проверка подстановкой возвращает правую часть уравнения.
 \item В № 10 скорость совместной работы \(\frac15+\frac1{12}=\frac{17}{60}\) резервуара в час; за два часа заполнится \(\frac{17}{30}\).
\end{itemize}
\end{document}
